\section{Leveled FHE from LWE}\label{fhe}

\subsection{The main idea: approximate eigenvalues}

Now we want to turn the public-key encryption from {[}lwe-crypto{]} into
a leveled FHE scheme. In other words: We want to be able to encrypt bits
(0s and 1s) and operate on them with AND and NOT gates.

It might help to imagine that, instead of AND and NOT, the operations we
want to encrypt are addition and multiplication. If \(x\) and \(y\) are
bits, then NOT \(x\) is just \(1 - x\), and \(x\) AND \(y\) is just
\(xy\). But it is easier to do algebra with \(+\) and \(\times\).

Recall the setup from {[}lwe-crypto{]}: We are going pick some large
integer \(q\) (in practice \(q\) could be anywhere from a few thousand
to \(2^{1000}\)), and do ``approximate linear algebra'' modulo \(q\). In
other words, we will do linear algebra, where all our calculations are
done modulo \(q\) -- but we will also allow the calculations to have a
small ``error'' \(\epsilon\), which will typically be much, much smaller
than \(q\).

As before, our \emph{secret key} will be a vector of length \(n\):
\[\mathbf{v} = \left( v_{1},\ldots,v_{n} \right) \in ({\mathbb{Z}}/q{\mathbb{Z}})^{n}.\]
Suppose we want to encode a message \(\mu\) that is just a single bit,
let us say \(\mu \in \left\{ 0,1 \right\}\). Our ciphertext will be a
square \(n\)-by-\(n\) matrix \(C\) such that
\[C\mathbf{v} \approx \mu\mathbf{v}.\] Now if we assume that
\(\mathbf{v}\) has at least one ``big'' entry (say \(v_{i}\)), then
decryption is easy: just compute the \(i\)-th entry of \(C\mathbf{v}\),
and determine whether it is closer to \(0\) or to \(v_{i}\).

strong(Remark):\\
With a bit of effort, it is possible to make this into a public-key
cryptosystem too. Just like in {[}lwe-crypto{]}, the main idea is to
release a table of vectors \(\mathbf{x}\) such that
\[\mathbf{x} \cdot \mathbf{v} \approx 0,\] and use that as a public key.
Given \(\mu\) and the public key, you can find a matrix \(C_{0}\) such
that \[C_{0}\mathbf{v} \approx 0\] then take
\[C = C_{0} + \mu\mathrm{\text{Id}},\] where \(\mathrm{\text{Id}}\) is
the identity matrix. This gives a \(C\) such that

\[C\mathbf{v} \approx \mu\mathbf{v}.\]

strong(Problem):\\
How do we build such a \(C_{0}\)? (One possible direction is to build it
row-by-row.)

\subsection{Operations on encrypted data}

To make homomorphic encryption work, we need to explain how to operate
on \(\mu\). We will describe three operations: addition, NOT, and
multiplication (aka AND).

Addition is simple: just add the matrices. If
\(C_{1}\mathbf{v} \approx \mu_{1}\mathbf{v}\) and
\(C_{2}\mathbf{v} \approx \mu_{2}\mathbf{v}\), then
\[\left( C_{1} + C_{2} \right)\mathbf{v} = C_{1}\mathbf{v} + C_{2}\mathbf{v} \approx \mu_{1}\mathbf{v} + \mu_{2}\mathbf{v} = \left( \mu_{1} + \mu_{2} \right)\mathbf{v}.\]
Of course, addition on bits is not a great operation, because if you add
\(1 + 1\), you get \(2\), and \(2\) is not a legitimate bit anymore. So
we will not really use this.

Negation of a bit (NOT) is equally simple. If
\(\mu \in \left\{ 0,1 \right\}\) is a bit, then its negation is simply
\(1 - \mu\). And if \(C\) is a ciphertext for \(\mu\), then
\(\mathrm{\text{Id}} - C\) is a ciphertext for \(1 - \mu\), since
\[\left( \mathrm{\text{Id}} - C \right)\mathbf{v} = \mathbf{v} - C\mathbf{v} \approx (1 - \mu)\mathbf{v}.\]

Multiplication is also a good operation on bits -- it is just AND. To
multiply two bits, you just multiply (matrix multiplication) the
ciphertexts:
\[C_{1}C_{2}\mathbf{v} \approx C_{1}\left( \mu_{2}\mathbf{v} \right) = \mu_{2}C_{1}\mathbf{v} \approx \mu_{2}\mu_{1}\mathbf{v} = \mu_{1}\mu_{2}\mathbf{v}.\]

At this point you might be concerned about this symbol \(\approx\) and
what happens to the size of the error. That is an important issue, and
we will resolve it with the help of a special operation called
``Flatten.''

Anyway, once you have AND and NOT, you can build arbitrary logic gates
-- and this is what we mean when we say you can perform arbitrary
calculations on your encrypted bits, without ever learning what those
bits are. At the end of the calculation, you can send the resulting
ciphertexts back to be decrypted.

\subsection{The ``Flatten''
operation}\label{a-constraint-on-the-secret-key-mathbfv-and-the-flatten-operation}

In order to make the error estimates work out, we are going to need to
make it so that all the ciphertext matrices \(C\) have ``small''
entries. In fact, we will be able to make it so that all entries of
\(C\) are either \(0\) or \(1\).

To make this work, we will assume our secret key \(\mathbf{v}\) has the
special form \[\mathbf{v} = \left. \begin{aligned}
( & a_{1},2a_{1},4a_{1},\ldots,2^{k}a_{1}, \\
 & a_{2},2a_{2},4a_{2},\ldots,2^{k}a_{2}, \\
 & \vdots \\
 & a_{r},2a_{r},4a_{r},\ldots,2^{k}a_{r})
\end{aligned} \right.,\] \phantomsection\label{fhe-v-form}{}

where \(k = \lfloor\log_{2}q\rfloor\).

To see how this helps us, try the following puzzle. Assume \(q = 11\)
(so all our vectors have entries modulo 11), and \(r = 1\), so our
secret key has the form
\[\mathbf{v} = \left( a_{1},2a_{1},4a_{1},8a_{1} \right).\] You know
\(\mathbf{v}\) has this form, but you do not know the specific value of
\(a_{1}\).

Now suppose we give you the vector \[\mathbf{x} = (9,0,0,0).\] We ask
you for another vector
\[\text{ Flatten}\left( \mathbf{x} \right) = \mathbf{x}',\] where
\(\mathbf{x}'\) has to have the following two properties:

\begin{itemize}
\item
  \(\mathbf{x}' \cdot \mathbf{v} = \mathbf{x} \cdot \mathbf{v}\), and
\item
  All the entries of \(\mathbf{x}'\) are either 0 or 1.
\end{itemize}

And you have to find this vector \(\mathbf{x}'\) without knowing
\(a_{1}\).

The solution is to use binary expansion: take
\(\mathbf{x}' = (1,0,0,1)\). You should check for yourself to see why
this works -- it boils down to the fact that \((1,0,0,1)\) is the binary
expansion of \(9\).

strong(Problem):\\
How would you flatten a different vector, like
\[\mathbf{x} = (9,3,1,4)?\] As a hint, remember we are working with
numbers modulo 11: So if you come across a number that is bigger than 11
in your calculation, it is safe to reduce it mod 11.

In general, if you know that \(\mathbf{v}\) has the form in
\hyperref[fhe-v-form]{{[}fhe-v-form{]}} and you are given some matrix
\(C\) with coefficients in \({\mathbb{Z}}/q{\mathbb{Z}}\), then you can
compute another matrix \(\text{Flatten}(C)\) such that:

\begin{itemize}
\item
  \(\text{Flatten}(C)\mathbf{v} = C\mathbf{v}\), and
\item
  All the entries of \(\text{Flatten}(C)\) are either 0 or 1.
\end{itemize}

The \(\text{Flatten}\) process is essentially the same binary-expansion
process we used above to turn \(\mathbf{x}\) into \(\mathbf{x}'\),
applied to each \(k + 1\) entries of each row of the matrix \(C\).

So now, using this \(\text{Flatten}\) operation, we can insist that all
of our ciphertexts \(C\) are matrices with coefficients in
\(\left\{ 0,1 \right\}\). For example, to multiply two messages
\(\mu_{1}\) and \(\mu_{2}\), we first multiply the corresponding
ciphertexts, then flatten the resulting product:
\[\text{ Flatten}\left( C_{1}C_{2} \right).\]

Of course, revealing that the secret key \(\mathbf{v}\) has this special
form will degrade security. This cryptosystem is as secure as an LWE
problem on vectors of length \(r\), not \(n\). So we need to make \(n\)
bigger, say \(n \approx r\log q\), to get the same level of security.

\subsection{Error analysis}\label{error-analysis}

Now let us compute more carefully what happens to the error when we add,
negate, and multiply bits. Suppose
\[C_{1}\mathbf{v} = \mu_{1}\mathbf{v} + \mathbf{\varepsilon}_{1},\]
where \(\mathbf{\varepsilon}_{1}\) is some vector with all its entries
bounded by some \(B\). (And similarly for \(C_{2}\) and \(\mu_{2}\).)

When we add two ciphertexts, the errors add:
\[\left( C_{1} + C_{2} \right)\mathbf{v} = \left( \mu_{1} + \mu_{2} \right)\mathbf{v} + \left( \mathbf{\varepsilon}_{1} + \mathbf{\varepsilon}_{2} \right).\]
So the error on the sum will be bounded by \(2B\).

Negation is similar to addition -- in fact, the error will not change at
all.

Multiplication is more complicated, and this is why we insisted that all
ciphertexts have entries in \(\left\{ 0,1 \right\}\). We compute
\[C_{1}C_{2}\mathbf{v} = C_{1}\left( \mu_{2}\mathbf{v} + \mathbf{\varepsilon}_{2} \right) = \mu_{1}\mu_{2}\mathbf{v} + \left( \mu_{2}\mathbf{\varepsilon}_{1} + C_{1}\mathbf{\varepsilon}_{2} \right).\]

Now since \(\mu_{2}\) is either \(0\) or \(1\), we know that
\(\mu_{2}\mathbf{\varepsilon}_{1}\) is a vector with all entries bounded
by \(B\). What about \(C_{1}\mathbf{\varepsilon}_{2}\)? Here we have to
think carefully about matrix multiplication: When you multiply an
\(n\)-by-\(n\) matrix by a vector, each entry of the product comes as a
sum of \(n\) different products. Now we are assuming that \(C_{1}\) is a
\(0\)-\(1\) matrix, and all entries of \(\mathbf{\varepsilon}_{2}\) are
bounded by \(B\)\ldots{} so the product has all entries bounded by
\(nB\). Adding this to the error for
\(\mu_{2}\mathbf{\varepsilon}_{1}\), we get that the total error in the
product \(C_{1}C_{2}\mathbf{v}\) is bounded by \((n + 1)B\).

In summary: We can start with ciphertexts having a very small error (if
you think carefully about this protocol, you will see that the error is
bounded by approximately \(n\log q\)). Every addition operation will
double the error bound; every multiplication (AND gate) will multiply it
by \((n + 1)\). And you cannot allow the error to exceed \(q/2\) --
otherwise the message cannot be decrypted. So you can perform
calculations of up to approximately \(\log_{n}q\) steps. (In fact, it is
a question of \emph{circuit depth}: You can start with many more than
\(\log_{n}q\) input bits, but no bit can follow a path of length greater
than \(\log_{n}q\) AND gates.)

This gives us a \emph{leveled} FHE protocol: It lets us evaluate
arbitrary circuits on encrypted data, as long as those circuits have
bounded depth. If we need to evaluate a bigger circuit, we have two
options:

\begin{enumerate}
\item
  Increase the value of \(q\). Of course, the cost of the computations
  increases with \(q\).
\item
  Use some technique to ``reset'' the error and start anew, as if with a
  freshly encrypted ciphertext. This approach is called
  \emph{bootstrapping} and it incurs some hefty computational costs. But
  for large circuits, it is the only viable option. Bootstrapping is
  beyond the scope of this book.
\end{enumerate}
